\documentclass[12pt,a4paper]{article}

% ============================================================
% PACKAGES
% ============================================================
\usepackage[utf8]{inputenc}
\usepackage[T1]{fontenc}
\usepackage[french]{babel}
\usepackage{geometry}
\usepackage{graphicx}
\usepackage{xcolor}
\usepackage{listings}
\usepackage{booktabs}
\usepackage{array}
\usepackage{hyperref}
\usepackage{fancyhdr}
\usepackage{titlesec}
\usepackage{enumitem}
\usepackage{mdframed}
\usepackage{tikz}
\usepackage{float}
\usepackage[table]{xcolor}

% ============================================================
% MISE EN PAGE
% ============================================================
\geometry{
  top=2.5cm,
  bottom=2.5cm,
  left=2.5cm,
  right=2.5cm
}

% ============================================================
% COULEURS
% ============================================================
\definecolor{primary}{RGB}{30, 100, 180}
\definecolor{secondary}{RGB}{240, 245, 255}
\definecolor{codebg}{RGB}{245, 245, 245}
\definecolor{codeframe}{RGB}{200, 200, 200}
\definecolor{success}{RGB}{39, 174, 96}
\definecolor{warning}{RGB}{230, 126, 34}
\definecolor{darkgray}{RGB}{60, 60, 60}

% ============================================================
% EN-TETE ET PIED DE PAGE
% ============================================================
\pagestyle{fancy}
\fancyhf{}
\fancyhead[L]{\small\textcolor{primary}{\textbf{Projet M1 GIL -- Gestion de Flotte}}}
\fancyhead[R]{\small\textcolor{darkgray}{Universite de Rouen | 2025--2026}}
\fancyfoot[C]{\thepage}
\renewcommand{\headrulewidth}{0.4pt}

% ============================================================
% STYLE DES TITRES
% ============================================================
\titleformat{\section}
  {\Large\bfseries\color{primary}}
  {\thesection.}{0.8em}{}[\titlerule]

\titleformat{\subsection}
  {\large\bfseries\color{darkgray}}
  {\thesubsection.}{0.6em}{}

\titleformat{\subsubsection}
  {\normalsize\bfseries\color{darkgray}}
  {\thesubsubsection.}{0.5em}{}

% ============================================================
% STYLE DU CODE
% ============================================================
\lstset{
  backgroundcolor=\color{codebg},
  frame=single,
  rulecolor=\color{codeframe},
  basicstyle=\ttfamily\small,
  keywordstyle=\color{primary}\bfseries,
  commentstyle=\color{gray}\itshape,
  stringstyle=\color{success},
  breaklines=true,
  showstringspaces=false,
  numbers=left,
  numberstyle=\tiny\color{gray},
  numbersep=8pt,
  tabsize=2,
  captionpos=b
}

% ============================================================
% BOITE INFO
% ============================================================
\newmdenv[
  backgroundcolor=secondary,
  linecolor=primary,
  linewidth=1.5pt,
  leftline=true,
  rightline=false,
  topline=false,
  bottomline=false,
  innerleftmargin=10pt,
  innerrightmargin=10pt,
  innertopmargin=8pt,
  innerbottommargin=8pt
]{infobox}

% ============================================================
% DOCUMENT
% ============================================================
\begin{document}

% ------------------------------------------------------------
% PAGE DE TITRE
% ------------------------------------------------------------
\begin{titlepage}
  \centering
  \vspace*{1cm}

  {\Large\textbf{Universite de Rouen Normandie}\\[0.3cm]
  \large Master 1 -- Genie Informatique et Logiciel}

  \vspace{1.5cm}
  \rule{\linewidth}{0.5mm}\\[0.4cm]

  {\Huge\bfseries\color{primary} Projet Gestion de Flotte de vehicules\\[0.3cm]}
  {\LARGE\bfseries Rapport de Mise en \oe{}uvre}\\[0.3cm]
  {\large\color{darkgray} Microservice Maintenance}

  \rule{\linewidth}{0.5mm}\\[1.5cm]

  \begin{minipage}{0.45\textwidth}
    \begin{flushleft}
      \textbf{Etudiant :}\\
      Florian Pepin \\[0.5cm]
      \textbf{Encadrants :}\\
      Lydia \& Luc
    \end{flushleft}
  \end{minipage}
  \hfill
  \begin{minipage}{0.45\textwidth}
    \begin{flushright}
      \textbf{Annee universitaire :}\\
      2025 -- 2026
    \end{flushright}
  \end{minipage}

  \vfill
  {\large\today}
\end{titlepage}

\tableofcontents
\newpage

% ------------------------------------------------------------
% INTRODUCTION
% ------------------------------------------------------------
\section{Objectifs du service Maintenance}

\begin{infobox}
\textbf{Objectif :} Implémenter le microservice de maintenance permettant la planification des interventions, le suivi de l'historique et la génération d'alertes préventives.
\end{infobox}

Le service Maintenance est responsable de la santé technique de la flotte. Ses fonctionnalités clés incluent :
\begin{itemize}
    \item \textbf{Planification} : Création et suivi des rendez-vous de maintenance.
    \item \textbf{Historique} : Conservation de toutes les interventions passées pour chaque véhicule.
    \item \textbf{Alertes Préventives} : Détection automatique des retards et des échéances kilométriques.
\end{itemize}

% ------------------------------------------------------------
% MODELISATION
% ------------------------------------------------------------
\section{Modélisation et Persistance}

Le service respecte le principe d'isolation des données : \textbf{un microservice = une base de données}. Il s'appuie sur une instance PostgreSQL dédiée (\texttt{postgres-maintenance}), dont le schéma est défini dans \texttt{db/03-maintenance.sql}.

\subsection{Base de données dédiée}
Contrairement aux autres services partageant l'instance principale, la maintenance dispose de sa propre base \texttt{maintenance\_db} tournant sur un conteneur isolé. Cette architecture garantit qu'une corruption ou une maintenance sur la base des véhicules n'impacte pas le suivi des entretiens.

\subsection{Entité MaintenanceRecord}
L'entité centrale \texttt{MaintenanceRecord} gère les cycles de vie des interventions via plusieurs énumérations (\texttt{Type}, \texttt{Status}, \texttt{Priority}).

\begin{lstlisting}[language=Java, caption=Entité MaintenanceRecord]
@Entity
@Table(name = "maintenance_records")
public class MaintenanceRecord {
    @Id
    @GeneratedValue(strategy = GenerationType.UUID)
    private UUID id;
    
    private UUID vehicleId; // Lien logique vers le Service Vehicules
    
    @Enumerated(EnumType.STRING)
    private MaintenanceStatus status = MaintenanceStatus.SCHEDULED;
    
    private LocalDate scheduledDate;
    private Integer nextServiceKm; // Seuil pour les alertes preventives
}
\end{lstlisting}

% ------------------------------------------------------------
% LOGIQUE METIER
% ------------------------------------------------------------
\section{Logique Métier et Alertes}

Le service implémente des mécanismes proactifs pour assurer le suivi de la maintenance.

\subsection{Gestion des retards (Overdue)}
Une tâche planifiée (\texttt{@Scheduled}) s'exécute quotidiennement pour identifier les maintenances prévues dont la date est dépassée.

L'expression Cron utilisée est \texttt{"0 0 1 * * ?"}. Elle se décompose comme suit :
\begin{itemize}
    \item \texttt{0} (1er) : 0 seconde.
    \item \texttt{0} (2e) : 0 minute.
    \item \texttt{1} (3e) : 1 heure du matin.
    \item \texttt{*} (4e) : Chaque jour du mois.
    \item \texttt{*} (5e) : Chaque mois.
    \item \texttt{?} (6e) : N'importe quel jour de la semaine.
\end{itemize}
Ainsi, le système vérifie automatiquement les retards chaque nuit à 01h00.

\begin{lstlisting}[language=Java, caption=Tâche de vérification des retards]
@Scheduled(cron = "0 0 1 * * ?")
@Transactional
public void checkOverdueMaintenance() {
    List<MaintenanceRecord> scheduledRecords = repository.findByStatus(MaintenanceStatus.SCHEDULED);
    LocalDate today = LocalDate.now();
    for (MaintenanceRecord record : scheduledRecords) {
        if (record.getScheduledDate().isBefore(today)) {
            record.setStatus(MaintenanceStatus.OVERDUE);
            repository.save(record);
            eventProducer.publishAlert(AlertEvent.overdue(
                record.getVehicleId(), "Maintenance en retard"));
        }
    }
}
\end{lstlisting}

% ------------------------------------------------------------
% INTEGRATION
% ------------------------------------------------------------
\section{Intégration et Communication}

\subsection{Architecture Événementielle}
Le service produit des messages sur le topic Kafka \texttt{alerts-topic}. Ces événements (\texttt{AlertEvent}) sont consommés par le service d'événements pour être notifiés aux gestionnaires.

\subsection{Exposition REST}
Une API complète a été développée pour permettre au Gateway de gérer les interventions :
\begin{center}
\begin{tabular}{|l|l|l|}
  \hline
  \rowcolor{primary!10}
  \textbf{Méthode} & \textbf{Endpoint} & \textbf{Action} \\
  \hline
  POST   & \texttt{/api/maintenance} & Planifier une intervention \\
  GET    & \texttt{/api/maintenance/vehicle/\{id\}} & Historique par véhicule \\
  PATCH  & \texttt{/api/maintenance/\{id\}/status} & Mettre à jour le statut \\
  \hline
\end{tabular}
\end{center}

\subsection{Exposition via le Gateway GraphQL}
Le service est pleinement intégré à l'API Gateway via GraphQL.

\subsubsection{Agrégation multi-services}
L'agrégation permet de croiser les données de trois services distincts :
\begin{itemize}
    \item \textbf{Maintenance $\rightarrow$ Véhicule} : Récupération des caractéristiques techniques (Service Véhicules).
    \item \textbf{Maintenance $\rightarrow$ Technicien} : Résolution de l'objet \texttt{Driver} lié au \texttt{technician\_id} (Service Conducteurs).
    \item \textbf{Véhicule $\rightarrow$ Maintenance} : Affichage de l'historique depuis la fiche véhicule.
\end{itemize}

\begin{lstlisting}[language=SQL, caption=Cohérence des types SQL/GraphQL]
-- Le champ JSONB parts_used en SQL est expose comme [String!]! en GraphQL
CREATE TABLE maintenance_records (
  ...
  parts_used JSONB NOT NULL DEFAULT '[]',
  ...
);
\end{lstlisting}

% ------------------------------------------------------------
% SAGA PATTERN
% ------------------------------------------------------------
\section{Architecture Événementielle (Saga Pattern)}

Le service participe au cycle de vie global via Kafka :
\begin{itemize}
    \item \textbf{Producers} : Publication sur \texttt{maintenance-topic} (\texttt{scheduled}, \texttt{started}, \texttt{completed}, \texttt{cancelled}).
    \item \textbf{Consumers} : Écoute du topic \texttt{vehicle-events} pour réagir aux mises à jour de kilométrage.
\end{itemize}

% ------------------------------------------------------------
% TESTS
% ------------------------------------------------------------
\section{Tests automatisés et couverture de code}

\begin{infobox}
\textbf{Objectif :} Garantir la robustesse via une couverture minimale de 80\%.
\end{infobox}

\subsection{Stratégie et Résultats}
\begin{itemize}
    \item \textbf{Tests Unitaires} : JUnit 5 + Mockito pour isoler la logique métier.
    \item \textbf{Tests d'Intégration} : MockMvc + base H2 pour valider l'API et le RBAC.
    \item \textbf{Couverture JaCoCo} : Supérieure à 80\% sur l'ensemble du périmètre métier.
\end{itemize}

% ------------------------------------------------------------
% DEPLOIEMENT
% ------------------------------------------------------------
\section{Déploiement Kubernetes et Helm}

\subsection{Isolation de l'infrastructure (Helm)}
\begin{itemize}
    \item \textbf{Chart fleet-infra} : Alias pour déployer un second StatefulSet (\texttt{postgres-maintenance-service}).
    \item \textbf{Initialisation} : ConfigMap \texttt{maintenance-db-init} injectant \texttt{03-maintenance.sql}.
\end{itemize}

\end{document}
